% ==============================================================================
% sec:zhodnoceni — souhrnné zhodnocení formou SWOT + verifikace hypotézy
% ==============================================================================

Předchozí analytické části poskytly komplexní empirický a~komparativní podklad pro~posouzení stavu sociálního dialogu v~\acloc{ČR}. K~syntéze těchto poznatků a~přechodu k~návrhové části slouží strukturovaná SWOT analýza, která hodnotí klíčové silné a~slabé stránky tuzemského systému, doplňuje je o~výhledové příležitosti a~identifikuje hlavní rizika.\par

\paragraph{Silné stránky (Strengths)}\label{par:swot_s}
\begin{itemize}[noitemsep]
	\item Funkční tripartitní platforma \acs{RHSD} s~ustáleným jednacím řádem, kontinuitou od~roku 1991 a~pevným ukotvením v~zákoně o~\acloc{KV} (č.~2/1991~Sb.) i~v~zákoníku práce.
	\item Uznané reprezentativní zastřešující organizace sociálních partnerů na~obou stranách trhu práce (\acs{ČMKOS}, \acs{ASO}, \acs{SPČR}, \acs{KZPS}).
	\item Institucionální existence právního nástroje rozšíření závaznosti \acpgen{KSVS} (\mbox{§\,7} zákona č.~2/1991~Sb.), umožňujícího plošné zvyšování smluvního pokrytí bez~nutnosti měnit primární legislativu.
	\item Vazba tuzemských dceřiných firem na~mateřské koncerny v~zemích s~rozvinutou demokratickou tradicí sociálního dialogu, zejména v~Německu a~Rakousku.
\end{itemize}

\paragraph{Slabé stránky (Weaknesses)}\label{par:swot_w}
\begin{itemize}[noitemsep]
	\item Extrémně nízká a~setrvale klesající odborová organizovanost v~\acloc{ČR} ($<$ \SI{12}{\percent}), což~zásadně omezuje vyjednávací legitimitu a~kapacitu zaměstnanecké strany.
	\item Nízké celkové pokrytí kolektivními smlouvami ($\approx$ \SI{30}{\percent}), což~je hluboko pod~unijním prahem \SI{80}{\percent} a~vyžaduje vypracování národního akčního plánu.\par
	\item Dlouhodobé podfinancování \acgen{APZ} ($\approx$ \SI{0{,}16}{\percent} \acs{HDP} oproti \SI{1{,}20}{\percent} v~Dánsku), což~limituje rozpočtový prostor pro rekvalifikace a~adaptaci na strukturální změny.
	\item Minimální využívání stávajícího mechanismu extenzí \acpgen{KSVS} v~praktické činnosti \acgen{MPSV}.
	\item Vyloučení nových a~atypických forem práce (ekonomicky závislé \acp{OSVČ}, švarcsystém, platformová práce) z~působnosti kolektivního vyjednávání.
	\item Asymetrický daňově-odvodový mix výrazně zatěžující standardní zaměstnaneckou práci, což motivuje obcházení zákoníku práce a~podkopává finanční udržitelnost sociálních systémů.
\end{itemize}

\paragraph{Příležitosti (Opportunities)}\label{par:swot_o}
\begin{itemize}[noitemsep]
	\item Transpoziční tlak evropské směrnice o~přiměřených minimálních mzdách a~směrnice o~transparentnosti odměňování, jež vyžadují posílení \acgen{KV} a~předložení národních akčních plánů.
	\item Možnost selektivní adaptace ghentského modelu svěřením odvětvové \acgen{APZ} a~rekvalifikačních programů reprezentativním odborovým svazům.
	\item Využití mezinárodních příkladů dobré praxe – zefektivnění rozšiřování \acpgen{KSVS} (vzor Rakousko) či zřízení sektorových fondů vzdělávání financovaných z~odvodů v~\acloc{KSVS} (vzor nizozemských fondů O\&O).
	\item Explicitní legislativní rozšíření působnosti \acpgen{KSVS} na samostatně činné osoby (solo-\acp{OSVČ}) se znaky ekonomické závislosti v~souladu s~Pokyny Evropské komise z~roku \SI{2022}{\rok}.\par
	\item Vybudování moderního digitálního registru kolektivních smluv k~odstranění informačních asymetrií na trhu práce.
\end{itemize}

\paragraph{Hrozby (Threats)}\label{par:swot_t}
\begin{itemize}[noitemsep]
	\item Uzamčení \acloc{ČR} v~pasti středních příjmů v~důsledku udržování subdodavatelského charakteru ekonomiky založeného na nízkých mzdách a~dlouhé pracovní době.
	\item Vyčleňování a~polarizace trhu práce pod vlivem stárnutí populace a~přílivu zahraniční pracovní síly bez~integračních mzdových standardů vyjednávaných sociálními partnery.
	\item Nekontrolovaný dopad automatizace a~generativní \acs{AI} v~podmínkách slabého odvětvového vyjednávání, což hrozí vytlačením ohrožených a~mladých pracovníků do prekarizovaných pozic bez~tarifní ochrany.
	\item Další emigrace tuzemského vysoce kvalifikovaného lidského kapitálu za lepšími mzdovými a~pracovními podmínkami do západních zemí s~koordinovaným sociálním dialogem.
	\item Eroze mezigenerační solidarity a~individualizace společenských rizik, kdy je nízký mzdový standard a~sociální podpojištění \acp{OSVČ} chápáno jako individuální selhání pracovníka.
	\item Politické tlaky na deregulaci a~pružnost trhu práce bez~odpovídajícího navýšení jistot, redukující sociální dialog na administrativní zátěž.
\end{itemize}

\paragraph{Verifikace hypotézy}\label{par:hypoteza_overeni}
Empirická zjištění této práce přinášejí konzistentní podklad pro vyhodnocení hypotézy formulované v~kap.~\ref{par:hypoteza}. Analýzy byly vedeny ve čtyřech komplementárních rovinách:
Deskriptivní rovina (kap.~\ref{ch:stav}) empiricky doložila institucionální slabost tuzemského vyjednávacího rámce, jenž vykazuje nízké a~klesající parametry pokrytí ($\approx$ \SI{31}{\percent}) i~členství (< \SI{12}{\percent}) souběžně s~podfinancovanou \acins{APZ}.\par
Komparativní rovina (kap.~\ref{sec:modely}) zařadila český model k~decentralizovaným systémům s~podnikovým těžištěm, přičemž mezinárodní srovnání ukázalo, že země s~obdobnou mzdovou hladinou a~strukturou (\ac{AT}) dosahují odlišných standardů díky rozvinuté mzdové koordinaci na odvětvové bázi.\par
Korelační rovina (kap.~\ref{sec:korelace}) prokázala silnou a~statisticky významnou asociaci mezi vysokým pokrytím \acpgen{KS} a~vyšší hodinovou produktivitou práce i~čistým hodinovým příjmem zaměstnanců, souběžně s~významnou zápornou korelací vůči délce týdenní pracovní doby a~jisté míry organizovanosti vůči rozptylu příjmů měřenému doloženým Giniho koeficientem.\par
Detailní případové studie (kap.~\ref{sec:case_studies}) potvrdily, že tam, kde odvětvové smlouvy (\acs{KSVS}) reálně působí, úspěšně zvyšují standardy odměňování a~rekvalifikací nad rámec legislativních minim, zatímco podnikové vyjednávání v~maloobchodě (Penny Market) bez~odvětvového zázemí vykazuje pouze minimální odchylku od zákonných prahů.\par

Souhrn těchto analytických rovin přináší odpovědi na vytyčené výzkumné otázky a~\textbf{podporuje hypotézu práce}: \emph{Posílení sociálního dialogu a~kolektivního vyjednávání v~\acloc{ČR} by přispělo k~udržitelnému hospodářskému rozvoji a~mzdové konvergenci}. Empiricky podložené inovace navrhované v~následující kapitole~\ref{ch:inovace} nepředstavují pouhá normativní doporučení, nýbrž jsou přímou reakcí na odhalené strukturální slabiny tuzemského trhu práce s~cílem uvolnit dosud nevyužitý koordinační potenciál sociálního dialogu.\par
